% ============================================================
%  第三章 曲面的第二基本形式
%  结构沿袭苏步青、胡和生《微分几何》第三章
%  记号约定：单位法向量 \mathbf{n}，第一/第二基本形式系数
%            E,F,G 与 L,M,N 及曲率 K,H,κ 均为标量不加粗。
% ============================================================

\section{第二基本形式}

设曲面的单位法向量为 $\mathbf{n}$（见式~\eqref{eq:曲面:单位法向量}）。曲面的第二基本形式描述曲面偏离其切平面的程度。

\begin{definition}[第二基本形式]
	曲面的第二基本形式定义为
	\begin{equation}
		\mathrm{II}=-\mathrm{d}\boldsymbol{r}\cdot\mathrm{d}\mathbf{n}
		=L\,\mathrm{d}u^{2}+2M\,\mathrm{d}u\,\mathrm{d}v+N\,\mathrm{d}v^{2},
		\label{eq:第二基本形式:定义}
	\end{equation}
	其中系数
	\begin{equation}
		L=\boldsymbol{r}_{uu}\cdot\mathbf{n},\qquad
		M=\boldsymbol{r}_{uv}\cdot\mathbf{n},\qquad
		N=\boldsymbol{r}_{vv}\cdot\mathbf{n}.
		\label{eq:第二基本形式:系数}
	\end{equation}
\end{definition}

\begin{remark}[几何意义]
	把第二基本形式限制在坐标曲线上可得 $L$、$N$ 的几何意义：在 $v$ 曲线上，$u$ 方向的法曲率由 $L/E$ 给出；在 $u$ 曲线上，$v$ 方向的法曲率由 $N/G$ 给出（法曲率的定义见下节）。因此 $L$、$N$ 反映了曲面沿两个坐标方向的弯曲程度。
\end{remark}

\section{曲面上曲线的曲率与法曲率}

设曲面上的曲线 $C$ 以弧长为参数：$\boldsymbol{r}=\boldsymbol{r}(u(s),v(s))$，其单位切向量为 $\mathbf{T}$，主法向量为 $\mathbf{N}$，曲率为 $\kappa$。由链式法则
\begin{equation}
	\boldsymbol{r}'=\boldsymbol{r}_u u'+\boldsymbol{r}_v v',
	\label{eq:曲面上曲线:切向量}
\end{equation}
\begin{equation}
	\boldsymbol{r}''=\boldsymbol{r}_{uu}(u')^{2}+2\boldsymbol{r}_{uv}u'v'+\boldsymbol{r}_{vv}(v')^{2}+\boldsymbol{r}_u u''+\boldsymbol{r}_v v''.
	\label{eq:曲面上曲线:二阶导数}
\end{equation}

\begin{definition}[法曲率]
	曲线 $C$ 在点处的曲率向量 $\kappa\mathbf{N}=\boldsymbol{r}''$ 在曲面法线方向上的分量称为该曲线沿方向 $u':v'$ 的\textbf{法曲率}：
	\begin{equation}
		\kappa_{n}=\boldsymbol{r}''\cdot\mathbf{n}
		=\frac{L\,\mathrm{d}u^{2}+2M\,\mathrm{d}u\,\mathrm{d}v+N\,\mathrm{d}v^{2}}
		{E\,\mathrm{d}u^{2}+2F\,\mathrm{d}u\,\mathrm{d}v+G\,\mathrm{d}v^{2}}
		=\frac{\mathrm{II}}{\mathrm{I}}.
		\label{eq:法曲率:定义}
	\end{equation}
\end{definition}

\begin{theorem}[Meusnier 定理]
	过曲面上一点、具有同一方向 $u':v'$ 的曲线在该点的法曲率都相同，且等于曲线主法向量 $\mathbf{N}$ 与曲面法线 $\mathbf{n}$ 的夹角 $\theta$ 满足
	\begin{equation}
		\kappa\cos\theta=\kappa_{n}.
		\label{eq:Meusnier:定理}
	\end{equation}
	特别地，取法截面（过 $\mathbf{n}$ 与所给方向的平面）截曲面所得的法截线，其曲率恰为 $\kappa_{n}$。
\end{theorem}

\section{主曲率与主方向}

\begin{definition}[主方向与主曲率]
	在曲面上一点，使法曲率取极值的方向称为\textbf{主方向}，相应的极值称为\textbf{主曲率}，记作 $\kappa_1,\kappa_2$。
\end{definition}

\begin{proposition}[主曲率方程]
	设 $\lambda=\mathrm{d}v/\mathrm{d}u$，由法曲率 $\kappa_n(\lambda)$ 的极值条件 $\mathrm{d}\kappa_n/\mathrm{d}\lambda=0$ 得
	\begin{equation}
		(M-\kappa F)+(N-\kappa G)\lambda=0.
		\label{eq:主方向:方程}
	\end{equation}
	由此导出的主曲率方程（特征方程）为
	\begin{equation}
		(EG-F^{2})\kappa^{2}-(EN-2FM+GL)\kappa+(LN-M^{2})=0.
		\label{eq:主曲率:特征方程}
	\end{equation}
	两个根 $\kappa_1,\kappa_2$ 即为主曲率；对应主方向满足
	\begin{equation}
		(L-\kappa_i E)\,\mathrm{d}u+(M-\kappa_i F)\,\mathrm{d}v=0,\qquad i=1,2.
		\label{eq:主方向:极值条件}
	\end{equation}
\end{proposition}

\begin{theorem}[Euler 公式]
	设 $\theta$ 是某方向与主方向 $\kappa_1$ 所对应方向之间的夹角（按第一基本形式的度量计），则该方向的法曲率为
	\begin{equation}
		\kappa_{n}(\theta)=\kappa_1\cos^{2}\theta+\kappa_2\sin^{2}\theta.
		\label{eq:Euler公式}
	\end{equation}
\end{theorem}

\begin{proposition}[Rodrigues 公式]
	沿主方向有
	\begin{equation}
		\mathrm{d}\mathbf{n}=-\kappa\,\mathrm{d}\boldsymbol{r},
		\label{eq:Rodrigues:公式}
	\end{equation}
	即沿主方向移动时，法向量的变化与位置向量的变化共线。这也是判别主方向的等价条件。
\end{proposition}

\section{高斯曲率与平均曲率}

\begin{definition}[高斯曲率与平均曲率]
	\textbf{高斯曲率} $K$ 与\textbf{平均曲率} $H$ 定义为
	\begin{equation}
		K=\kappa_1\kappa_2=\frac{LN-M^{2}}{EG-F^{2}},
		\label{eq:高斯曲率:定义}
	\end{equation}
	\begin{equation}
		H=\frac{\kappa_1+\kappa_2}{2}
		=\frac{EN-2FM+GL}{2(EG-F^{2})}.
		\label{eq:平均曲率:定义}
	\end{equation}
\end{definition}

\begin{theorem}[高斯绝妙定理（Theorema Egregium）]
	高斯曲率 $K$ 是内蕴量，它只由第一基本形式及其导数决定，与曲面在空间中的嵌入方式无关。
	\label{thm:高斯绝妙定理}
\end{theorem}

在正交参数曲线网（$F=0$）下，高斯曲率有简洁的显式公式：
\begin{equation}
	K=-\frac{1}{2\sqrt{EG}}\left[
		\frac{\partial}{\partial u}\left(\frac{G_u}{\sqrt{EG}}\right)
		+\frac{\partial}{\partial v}\left(\frac{E_v}{\sqrt{EG}}\right)
	\right].
	\label{eq:高斯曲率:正交网公式}
\end{equation}

\begin{example}[球面与圆柱面的高斯曲率]
	对单位球面取式~\eqref{eq:球面:参数表示}，有 $E=1$、$G=\cos^{2}u$、$F=0$，代入式~\eqref{eq:高斯曲率:正交网公式} 得 $K=1$。
	对圆柱面 $\boldsymbol{r}(u,v)=(a\cos v,a\sin v,u)$，有 $E=1$、$F=0$、$G=a^{2}$，故 $K=0$。由此可见 $K$ 在等距变换下不变，而平面与圆柱面等距（$K=0$），球面与平面不等距（$K\neq0$）。
\end{example}

\section{直纹面与可展曲面}

\begin{definition}[直纹面]
	由一族直线（直母线）构成的曲面称为\textbf{直纹面}，其参数表示为
	\begin{equation}
		\boldsymbol{r}(u,v)=\boldsymbol{a}(u)+v\,\boldsymbol{b}(u),
		\label{eq:直纹面:定义}
	\end{equation}
	其中 $\boldsymbol{a}(u)$ 为\textbf{准线}，$\boldsymbol{b}(u)\neq\boldsymbol{0}$ 为直母线方向。
\end{definition}

\begin{definition}[可展曲面]
	如果直纹面沿每条直母线有同一个切平面，即法向量沿直母线方向不变，则称为\textbf{可展曲面}。
\end{definition}

\begin{theorem}[可展曲面的判别]
	直纹面 $\boldsymbol{r}(u,v)=\boldsymbol{a}(u)+v\boldsymbol{b}(u)$ 是可展的，当且仅当
	\begin{equation}
		\bigl(\boldsymbol{a}'(u),\;\boldsymbol{b}(u),\;\boldsymbol{b}'(u)\bigr)=0,
		\label{eq:可展曲面:判别}
	\end{equation}
	即混合积恒为零；等价地，$K=0$（高斯曲率处处为零）。
\end{theorem}

\begin{theorem}[可展曲面的分类]
	除平面外，可展曲面必为下列三类之一：
	\begin{enumerate}
		\item \textbf{柱面}：$\boldsymbol{b}(u)$ 是常方向向量；
		\item \textbf{锥面}：$\boldsymbol{a}(u)$ 退化为一个定点（所有直母线共点）；
		\item \textbf{切线曲面}：$\boldsymbol{b}(u)\parallel\boldsymbol{a}'(u)$，即直母线为某一空间曲线的切线。
	\end{enumerate}
\end{theorem}

\begin{remark}
	可展曲面可以无皱褶地展开成平面，其本质原因是 $K=0$（高斯绝妙定理保证了这一几何性质，见式~\eqref{eq:高斯曲率:正交网公式}）。
\end{remark}

\section{脐点与极小曲面}

\begin{definition}[脐点]
	若在某点处主曲率相等（$\kappa_1=\kappa_2$），即该点的法曲率与方向无关，则称该点为\textbf{脐点}。$\kappa_1=\kappa_2=0$ 的脐点称为\textbf{平点}。
\end{definition}

\begin{theorem}[全脐点曲面]
	如果曲面上每一点都是脐点，则该曲面必为平面或球面（在充分光滑且连通的假设下）。
	\label{thm:全脐点曲面}
\end{theorem}

\begin{definition}[极小曲面]
	平均曲率处处为零（$H=0$）的曲面称为\textbf{极小曲面}。它的直观意义是：以给定边界为界的极小曲面使面积取驻值。
\end{definition}

\begin{example}[悬链面与正螺旋面]
	\textbf{悬链面}（catenoid）
	\begin{equation}
		\boldsymbol{r}(u,v)=\bigl(a\cosh u\cos v,\;a\cosh u\sin v,\;u\bigr)
		\label{eq:悬链面:参数表示}
	\end{equation}
	与\textbf{正螺旋面}（helicoid）
	\begin{equation}
		\boldsymbol{r}(u,v)=\bigl(u\cos v,\;u\sin v,\;v\bigr)
		\label{eq:正螺旋面:参数表示}
	\end{equation}
	都是极小曲面，且二者之间存在保形变换（它们互为共形等价，是经典的对偶例子）。
\end{example}

\section{曲面论基本定理}

首先引入第二类 Christoffel 符号 $\Gamma_{ij}^{k}$，把切平面的二阶导数按标架 $\{\boldsymbol{r}_u,\boldsymbol{r}_v,\mathbf{n}\}$ 分解（Gauss 方程）：
\begin{equation}
	\boldsymbol{r}_{uu}=\Gamma_{11}^{1}\boldsymbol{r}_u+\Gamma_{11}^{2}\boldsymbol{r}_v+L\mathbf{n},
	\qquad
	\boldsymbol{r}_{uv}=\Gamma_{12}^{1}\boldsymbol{r}_u+\Gamma_{12}^{2}\boldsymbol{r}_v+M\mathbf{n},
	\qquad
	\boldsymbol{r}_{vv}=\Gamma_{22}^{1}\boldsymbol{r}_u+\Gamma_{22}^{2}\boldsymbol{r}_v+N\mathbf{n}.
	\label{eq:Gauss方程:分解}
\end{equation}

\begin{theorem}[Gauss--Codazzi 方程]
	由式~\eqref{eq:Gauss方程:分解} 的可积性条件（混合偏导次序交换）可推出两类相容性方程：
	\textbf{Gauss 方程}断言高斯曲率 $K$ 由第一基本形式确定（即高斯绝妙定理，见式~\eqref{eq:高斯曲率:正交网公式}）；\textbf{Codazzi--Mainardi 方程}为
	\begin{equation}
		\frac{\partial L}{\partial v}-\frac{\partial M}{\partial u}
		=L\Gamma_{12}^{1}+M\bigl(\Gamma_{12}^{2}-\Gamma_{11}^{1}\bigr)-N\Gamma_{11}^{2},
		\label{eq:Codazzi:方程1}
	\end{equation}
	\begin{equation}
		\frac{\partial M}{\partial v}-\frac{\partial N}{\partial u}
		=L\Gamma_{22}^{1}+M\bigl(\Gamma_{22}^{2}-\Gamma_{12}^{1}\bigr)-N\Gamma_{12}^{2}.
		\label{eq:Codazzi:方程2}
	\end{equation}
\end{theorem}

\begin{theorem}[曲面论基本定理（Bonnet 定理）]
	设给定两个二次微分形式
	\begin{equation}
		\mathrm{I}=E\,\mathrm{d}u^{2}+2F\,\mathrm{d}u\,\mathrm{d}v+G\,\mathrm{d}v^{2},
		\qquad
		\mathrm{II}=L\,\mathrm{d}u^{2}+2M\,\mathrm{d}u\,\mathrm{d}v+N\,\mathrm{d}v^{2},
		\label{eq:Bonnet:基本形式}
	\end{equation}
	其中 $\mathrm{I}$ 正定，且系数满足 Gauss--Codazzi 方程，则存在正则曲面以 $\mathrm{I}$、$\mathrm{II}$ 为其第一、第二基本形式；并且这样的曲面在刚体运动的意义下\textbf{唯一}。
	\label{thm:Bonnet定理}
\end{theorem}

\begin{remark}
	第一基本形式决定曲面的度量（内蕴几何），第二基本形式决定曲面在空间中的弯曲方式（外在几何）。Bonnet 定理把两者"拼合"成曲面，其证明本质上是解一阶线性偏微分方程组的可积性（Frobenius 定理），与外微分形式方法（第四章）密切相关。
\end{remark}
